%===========================================================
\chapter{Predictive Modelling, Constrained Selection and Biological Interpretation}
\label{chap:model_training_evaluation}
%===========================================================

%-----------------------------------------------------------
\section{Introduction}
%-----------------------------------------------------------

Chapter~\ref{chap:feature_selection} established a statistically stable and 
structurally diversified CpG panel of fixed cardinality ($K = 5{,}000$) for 
each dataset. The proposed pipeline integrated empirical variability screening, 
biologically weighted stability selection, correlation-based redundancy pruning, 
and genomic diversification, ensuring that the resulting signatures are 
reproducible within each cohort and structurally constrained by design. 
The inter-dataset evaluation of Section~5.4 further revealed that cross-cohort 
replicability is partial and asymmetric: GSE69914 and GSE225845 share a 
concordant drift orientation and support bidirectional zero-shot transfer, 
whereas GSE287331 exhibits a systematically inverted drift axis, yielding 
sub-random transfer performance against both partner cohorts.

These structural properties define the design space for the analyses developed 
in the present chapter. The selected CpG panels are characterised by statistical 
robustness and biological diversification, but their capacity to discriminate 
Normal from Adjacent tissue under supervised modelling has not yet been 
assessed. Leakage-controlled feature selection guarantees that the chosen loci 
are statistically stable and non-redundant; it does not, by itself, imply that 
they encode sufficient discriminative information for a classifier to operate 
reliably on held-out data.

The objectives of this chapter are threefold. First, intra-dataset 
classification performance is quantified to verify whether the selected CpG 
panels support accurate discrimination of Normal and Adjacent tissue within 
each cohort. A linear Support Vector Machine is adopted as the primary 
classifier for GSE225845 and GSE287331, a $k$-Nearest Neighbours introduced for 
GSE69914 to assess the degree of linear separability. Second, cross-dataset 
transferability is examined by training on one cohort and evaluating directly 
on another, with no information from the target dataset entering the training 
procedure. This zero-leakage design quantifies the extent to which the inferred 
epigenetic drift axis generalises across independent cohorts, and whether 
performance degradation is systematically associated with the directional 
concordance structure documented in Chapter~\ref{chap:feature_selection}. 
Third, a constrained combinatorial compression step is introduced to derive a 
compact 50-CpG subset through a mixed-integer linear programming formulation 
inspired by the classical 0--1 knapsack paradigm. This step investigates 
whether predictive structure and biological diversity can be jointly preserved 
under severe dimensional reduction, subject to explicit constraints on 
cardinality, chromosomal distribution, correlation structure, and 
methylation directionality.

Finally, the biological coherence of the compressed panels is evaluated through 
CpG-to-gene mapping, cross-referenced against the COSMIC Cancer Gene Census (CGC, version 103, GRCh38)~\cite{sondka2018cosmic,cosmic_cgc_v103}, 
and pathway-level enrichment analysis. This interpretive step closes the loop 
between statistical discrimination and functional relevance, assessing whether 
the loci selected by a purely data-driven optimisation converge on genomic 
regions with established roles in cancer-related epigenetic deregulation.

Together, these analyses progress from statistically robust feature selection 
to supervised evaluation, external validation, and structured biological 
interpretation. The limitations that emerge — in particular the failure of 
linear $\Delta\beta$ signatures to generalise across cohorts with divergent 
drift orientations — directly motivate the representation-learning approach 
developed in Chapter~\ref{chap:CpGPT}.

%-----------------------------------------------------------
\section{Predictive Framework}
%-----------------------------------------------------------

%-----------------------------------------------------------
\subsection{Intra-Dataset Supervised Modelling}
%-----------------------------------------------------------

This section describes the supervised classification framework
applied independently within each cohort.

\paragraph{Data partitioning and leakage control}
Description of train/test split strategy and preprocessing
performed strictly on training data.

\paragraph{Linear Support Vector Machine}
Formal definition of the linear SVM classifier,
regularisation parameter selection,
and decision function.

\paragraph{k-Nearest Neighbours (GSE69914 only)}
Definition of the kNN classifier,
distance metric,
and rationale for inclusion as a local non-parametric baseline.

\paragraph{Performance metrics}
Definition of AUC, balanced accuracy,
Matthews correlation coefficient,
precision, recall and F1-score.

%-----------------------------------------------------------
\subsection{Cross-Dataset Evaluation Protocol}
%-----------------------------------------------------------

This section defines the cross-cohort transferability framework.

\paragraph{Feature space alignment}
Restriction to shared CpG space and column ordering consistency.

\paragraph{Train-on-A / Test-on-B design}
Formal description of the external validation setting,
including absence of retraining on the target cohort.

\paragraph{Evaluation criteria}
Definition of discrimination metrics and interpretation
under distributional shift.

%-----------------------------------------------------------
\subsection{Constrained Subset Optimisation: Knapsack Formulation}
%-----------------------------------------------------------

This section formalises the compression of the 5,000-CpG panel
into a fixed-cardinality subset.

\paragraph{Objective function}
Definition of the optimisation criterion combining
statistical ranking and diversification score.

\paragraph{Cardinality and structural constraints}
Fixed subset size, chromosome caps,
genomic context constraints,
correlation-exclusion rules,
and sign-balance conditions.

\paragraph{Optimisation strategy}
Description of greedy or mixed-integer approximation scheme
and computational considerations.

%-----------------------------------------------------------
\subsection{Gene Set and Functional Enrichment Analysis}
%-----------------------------------------------------------

This section evaluates the biological coherence
of the selected CpG panels.

\paragraph{CpG-to-gene mapping}
Annotation procedure based on Illumina manifest.

\paragraph{Enrichment testing}
Definition of over-representation analysis,
background universe,
multiple-testing correction.

\paragraph{Reference databases}
Description of external resources
(e.g. COSMIC Cancer Gene Census, curated field-cancerization genes).

%-----------------------------------------------------------
\section{Intra-Dataset Results}
%-----------------------------------------------------------

%-----------------------------------------------------------
\subsection{Classification Performance on the 5,000-CpG Panel}
%-----------------------------------------------------------

Performance metrics obtained independently within each cohort.

%-----------------------------------------------------------
\subsection{Comparison Between Linear and Local Classifiers}
%-----------------------------------------------------------

Comparison between SVM and kNN (GSE69914),
with analysis of linear separability.

%-----------------------------------------------------------
\subsection{Compression to a 50-CpG Panel}
%-----------------------------------------------------------

Evaluation of predictive stability after knapsack-based reduction.

%-----------------------------------------------------------
\subsection{Biological Characterisation of the 50-CpG Signature}
%-----------------------------------------------------------

Functional enrichment results and gene-level interpretation.

%-----------------------------------------------------------
\section{Inter-Dataset Results}
%-----------------------------------------------------------

%-----------------------------------------------------------
\subsection{Cross-Cohort Transferability of the 5,000-CpG Panel}
%-----------------------------------------------------------

External validation results across all train/test combinations.

%-----------------------------------------------------------
\subsection{Directional Concordance and Performance Degradation}
%-----------------------------------------------------------

Analysis of performance variation in relation to
inter-dataset $\Delta\beta$ concordance.

%-----------------------------------------------------------
\subsection{Cross-Dataset Behaviour of the 50-CpG Subset}
%-----------------------------------------------------------

Evaluation of compressed signature under distributional shift.

%-----------------------------------------------------------
\section{Synthesis and Implications}
%-----------------------------------------------------------

This section summarises predictive findings,
identifies structural limitations of linear
$\Delta\beta$-based signatures,
and motivates the transition toward
representation-learning approaches
in the subsequent chapter.